From - Mon May 22 18:50:05 2000
Received: from imo-d05.mx.aol.com (imo-d05.mx.aol.com [205.188.157.37])
	by math.usu.edu (8.9.3+Sun/8.9.1) with SMTP id WAA24825
	for <symanzik@sunfs.math.usu.edu>; Sun, 21 May 2000 22:21:55 -0600 (MDT)
From: Case570@netscape.net
Received: from Case570@netscape.net
	by imo-d05.mx.aol.com (mail_out_v27.9.) id z.ec.bea9 (16228)
	 for <symanzik@sunfs.math.usu.edu>; Mon, 22 May 2000 00:23:40 -0400 (EDT)
Message-ID: <ec.bea9.265a104b@netscape.net>
Date: Mon, 22 May 2000 00:23:39 EDT
Subject:  Homework#6
To: symanzik@sunfs.math.usu.edu
MIME-Version: 1.0
Content-Type: text/plain; charset="US-ASCII"
Content-Transfer-Encoding: 7bit
X-Mailer: Unknown
Content-Length: 6054
Status: RO
X-Mozilla-Status: 8001
X-Mozilla-Status2: 00000000
X-UIDL: 38baa60a00000304

Electronic Textbook and Teachware Tools

Glen Erickson

Stats 5810 
Homework # 6


For this project we looked at and reviewed several teaching tools and textbooks and from there decided  exactly what one to focus on in our presentation. Following are a list of some of the texts I looked at and some of the good and bad points.
     Hyperstats: This text was interesting and had lots and lots of material in it, however most of the information was too advanced for beginning stats students. The links were good to some areas and poor for other areas. The section that covered Histograms,  stem and box plots was not really easy to understand and I felt that students would have a harder time understanding this concept if I were to teach from this text. I did like sone of the examples and practice exercises but felt like the large amount of reading would be to much for beginning students.
    
      Seeing Statistics: This site was fun and contained some useful information on a great deal of statistics but did not have a lot of or  very detailed text on histograms or box plots. I did not find many interactive parte that allowed students to participate in actually entering data and seeing how that different data changed the histogram or box plot. I did like many of the features of this site. The simplicity of the text was appealing to me and the site was easy to move around in but I thought it was incomplete in some areas. 
     UCLA Statistics Electronic Textbook: This text book was a little confusing to me and I got lost when I went to some of the links. There was one link that I went to and it would not let me back. Once I found the section on the information I was looking for I quite liked it but felt that it was to hard to get to and many of the beginners like myself would get lost to easy. It had some good features like the words that were important were in a different color and when you clicked on them a definition came up. I did not find many interactive sections and got bored so I went to another site.
     Sufrstat: This site I found to be very user friendly and easy to move around in. The text was fairly simple to read and was understandable. The interactive sections were a little weak and that was one of the things I was looking for. This site did have some good links and other texts as references to look at. I was looking for something fairly simple and easy to teach and up to this point had not found it.
    


      SticiGui: When I first saw the name I was intrigued and as I got into the site I liked it even more. The text was written in a easy to nice manner and was simple to understand. The examples were very useful and fun to read and the interactive parts seemed to put it all together for me. I especially liked the definitions that came up when ever you hit a colored work. I was able to find all the information that I looked for and did not get lost once, for me that is important! 
     Statsoft: Statsoft is a pretty good text because it is well written and pretty easy to understand. I tried to find some interactive exercises and did but they were not very fun nor did they show the student much about histograms.
     Cyberk.com: This site was very good and it did not take me long to decide to use this one for my presentation. This site offered more of what I was looking for than any other site. The section on histograms was easy to understand and fun to read . The subject was presented in an interesting and fun way. The interactive and graphic examples were great and I am a sucker for nice graphics. The practice examples were easy to find and work on and the solutions were in the back for reference. After using this site for a short time and trying some of the examples. I felt like I had a better grasp of the histogram and its use in statistics than from any other site. 
     From the several presentations given in class it became clear to me at least that a teaching presentation given solely from the net was not as successful as a presentation given partly on the board, and partly on the web. The combination of the two can produce an interesting and helpful setting for the learning and understanding of statistics. The traditional presentation of writing only on the board has now been rescued by the use of interactive teaching tools and texts on the world wide web. Using these two in a combination of traditional and new technology can liven up the old boring lecture and spark the interest of new and old students. One must be careful not to load to much of one or the other into a lecture. The traditional writing on the board still has its place and can be effective in teaching basic concepts. With the addition of some interactive web sites for use as examples and practice exercises a lecture becomes fun and students take some initiative to lear!
n on their own. This self motivation thing is hard to come by at any level but especially in students. I found myself coming home and playing around on the sites just to see what the different data would to the graphs and scatter plots and other graphics. The use of some web sites sparked an interest in me that I did not have previously.
I felt like the presentation that Bruce and I gave on histograms and box plots contained some good information and good references but we had to much web based material.
We should have used the board to describe and teach and used the web to show how these concepts worked in real life examples. The web can enhance and add some much needed color and spark to many lectures. I know that these tools and texts can be a great asset in teaching stats at any level.  
     I have not been able to get hold of Bruce to help write this part of the assignment so I am doing it myself as the due date is upon us. I will continue to try to reach him and see what he is going to do about the class. I will keep you informed when any new information becomes available.  
 

----------
Get your own FREE, personal Netscape Webmail account today at http://home.netscape.com/webmail/
